\section{Implementación del Software}\label{sec:implementation}

\subsection{Entorno tecnológico}
El desarrollo se realizó íntegramente en \emph{Google Colab}, sobre un cuaderno
Jupyter (formato \texttt{.ipynb}) ejecutado con un kernel de Python~3. El
paradigma de cómputo es exclusivamente \emph{CPU}; no se emplearon librerías de
aprendizaje profundo (PyTorch, TensorFlow) ni implementaciones de \emph{gradient
boosting} (XGBoost, LightGBM), dado que el clasificador seleccionado ---un
ensamble Random Forest--- está disponible de forma nativa en \emph{scikit-learn}.
Cabe señalar, como contraste, que el autor original del repositorio de datos
empleó un clasificador XGBoost sobre la versión filtrada del conjunto; el presente
trabajo constituye un pipeline independiente basado en scikit-learn.

La Tabla~\ref{tab:stack} resume el ecosistema de librerías identificado en las
celdas de importación del cuaderno.

\begin{table}[H]
\centering
\caption{Ecosistema tecnológico del proyecto}
\label{tab:stack}
\footnotesize
\renewcommand{\arraystretch}{1.2}
\begin{tabular}{p{2.2cm}p{5.4cm}}
\toprule
\textbf{Librería} & \textbf{Rol en el pipeline} \\
\midrule
\texttt{numpy} & Álgebra vectorial y control de semilla aleatoria \\
\texttt{pandas} & Estructura de datos tabular e ingesta \\
\texttt{matplotlib} & Visualización científica base \\
\texttt{seaborn} & Visualización estadística (densidades, heatmaps) \\
\texttt{kagglehub} & Ingesta programática del dataset desde Kaggle \\
\texttt{scikit-learn} & Preprocesamiento, PCA, modelado, evaluación \\
\bottomrule
\end{tabular}
\end{table}

Dentro de \emph{scikit-learn}, los módulos utilizados son:
\texttt{model\_selection} (\texttt{train\_test\_split}, \texttt{GridSearchCV}),
\texttt{pipeline} (\texttt{Pipeline}), \texttt{impute} (\texttt{SimpleImputer}),
\texttt{preprocessing} (\texttt{RobustScaler}), \texttt{compose}
(\texttt{ColumnTransformer}), \texttt{decomposition} (\texttt{PCA}),
\texttt{ensemble} (\texttt{RandomForestClassifier}) y \texttt{metrics}
(\texttt{classification\_report}, \texttt{confusion\_matrix},
\texttt{roc\_curve}, \texttt{auc}, \texttt{precision\_recall\_curve},
\texttt{average\_precision\_score}, \texttt{cohen\_kappa\_score},
\texttt{make\_scorer}, \texttt{f1\_score}).

\subsection{Organización modular del código}
El proyecto se implementó como un cuaderno lineal estructurado en bloques
funcionales cohesivos, en lugar de un paquete distribuido en un directorio
\texttt{/src}. Cada bloque encapsula una etapa del pipeline y expone sus
resultados a la siguiente mediante variables de estado explícitas. La
Tabla~\ref{tab:modules} describe esta organización modular por bloques.

\begin{table}[H]
\centering
\caption{Organización modular del cuaderno por bloques funcionales}
\label{tab:modules}
\footnotesize
\renewcommand{\arraystretch}{1.2}
\begin{tabular}{p{1.0cm}p{3.0cm}p{3.4cm}}
\toprule
\textbf{Bloque} & \textbf{Responsabilidad} & \textbf{Artefactos producidos} \\
\midrule
1 & Configuración global e ingesta & \texttt{df}, semilla global \\
2 & Diagnóstico y corrección de anomalías & \texttt{df} corregido, tasa de nulos \\
3 & Ingeniería de etiqueta y correlación & \texttt{Target\_Helada}, matriz Pearson \\
4 & Partición y preprocesamiento & \texttt{Pipeline}, \texttt{X\_train/test} escalados \\
5 & Reducción de dimensionalidad & \texttt{PCA}, \texttt{X\_train/test\_pca} \\
6 & Modelado y optimización & \texttt{best\_model}, hiperparámetros \\
7 & Evaluación integral & métricas, curvas ROC/PR, Kappa \\
\bottomrule
\end{tabular}
\end{table}

El preprocesamiento se encapsuló mediante los abstractores composicionales de
scikit-learn: un \texttt{Pipeline} que encadena \texttt{SimpleImputer} y
\texttt{RobustScaler}, integrado a su vez en un \texttt{ColumnTransformer}. Este
diseño garantiza que las transformaciones se ajusten (\emph{fit}) exclusivamente
sobre el conjunto de entrenamiento y se apliquen (\emph{transform}) de forma ciega
sobre el conjunto de prueba, previniendo la fuga de información estadística entre
particiones.

\subsection{Gestión de dependencias}
Las dependencias se gestionan a través del entorno preinstalado de Google Colab,
complementado por una instrucción de instalación explícita
(\texttt{pip install kagglehub[pandas-datasets] seaborn matplotlib scikit-learn})
documentada en la primera celda del cuaderno. El proyecto no fija versiones
específicas mediante un archivo \texttt{requirements.txt} o \texttt{environment.yml};
la adopción de un manifiesto de dependencias con versiones ancladas se plantea
como una mejora para robustecer la reproducibilidad a largo plazo.

\subsection{Reproducibilidad determinista}
El determinismo del pipeline se garantiza mediante la fijación de una semilla
global $\text{SEED}=42$, propagada por una función dedicada,
\texttt{set\_global\_seed}, que controla las tres fuentes de aleatoriedad activas
en el entorno de ejecución:
\begin{enumerate}
\item \texttt{random.seed(seed)} --- generador pseudoaleatorio de la biblioteca
estándar de Python.
\item \texttt{np.random.seed(seed)} --- generador de NumPy, base de las
operaciones estocásticas de scikit-learn.
\item \texttt{os.environ['PYTHONHASHSEED'] = str(seed)} --- fijación del hash de
Python para estabilizar el orden de estructuras dependientes de hashing.
\end{enumerate}

Esta misma semilla se reutiliza de forma consistente en cada componente
estocástico del pipeline: el muestreo de la matriz de correlación
(\texttt{sample(random\_state=SEED)}), la partición estratificada
(\texttt{train\_test\_split(random\_state=SEED)}), la descomposición espectral
(\texttt{PCA(random\_state=SEED)}) y la inicialización del ensamble
(\texttt{RandomForestClassifier(random\_state=SEED)}). Dado que el cómputo se
ejecuta enteramente sobre CPU mediante scikit-learn, no intervienen fuentes de no
determinismo asociadas a aceleradores por hardware (p.\,ej. operaciones no
deterministas de CUDA/cuDNN), por lo que no se requieren mecanismos adicionales de
control de aleatoriedad a nivel de GPU. La reproducibilidad experimental queda así
garantizada de extremo a extremo bajo el entorno descrito.